\section{Additional result figures and the repair gate}
\label{app:extrafigs}

% Moved from the main paper's §5 for the 3DV page limit (2026-07-13):
% fig:gen (S2), fig:tails (S3), fig:nsig (S4), and the stage-3a gate
% figure fig:toys (S5) with its mechanism prose. Captions and % src
% comments carried over verbatim; the verdicts these figures illustrate
% are stated with their numbers in main-paper §5.
% Order note (2026-07-13 later): nsig moved BEFORE toys — same-class
% floats place in source order, and the tall toys figure was stranding
% nsig alone on a fourth page. S4/S5 swapped accordingly (main-text
% plain-text refs updated in results.tex).

Fig.~\ref{fig:gen} shows the generality-wave training trajectories
(main paper \S5.2, Table~2); Fig.~\ref{fig:tails} the per-seed tail
spread behind the sphere row of main-paper Table~1;
% (round-4 final numbering: main Table 1 = C-matrix, Table 2 = generality)
Fig.~\ref{fig:nsig} the torus significant-count check behind
\S5.1's zero-phantom-handles claim.

\begin{figure*}[t]
\centering
% src: topo3/report/{bob,fandisk}_series.png (topo_loss_report.py)
\includegraphics[width=0.495\linewidth]{bob_series.png}\hfill
\includegraphics[width=0.495\linewidth]{fandisk_series.png}
\caption{Generality trajectories (line = seed mean, band = seed min--max).
\emph{Left}: bob (genus 1) --- the H1 result off the analytic family: the
loss halves the loop error with $\#$sig H1 $= 2$ in every seed, while the
control rides the baseline after collapsing a loop. \emph{Right}: fandisk
(CAD) --- the largest effect measured in this study ($10.4\times$).}
\label{fig:gen}
\end{figure*}

\begin{figure}[t]
\centering
% src: topo3/report/sphere_tail.png (per-seed tail means; results.json)
\includegraphics[width=0.9\linewidth]{sphere_tail.png}
\caption{Sphere, per-seed tail bottlenecks (dots) with arm means (bars).
Every loss arm sits far below baseline with tight seed spread; the two
control strengths bracket baseline from \emph{above}; the $\rho$ ablation
arms are indistinguishable from $\rho{=}0.1$.}
\label{fig:tails}
\end{figure}

\begin{figure}[t]
\centering
% src: topo3/report/torus_nsig.png (seed-mean significant-H1 count; axis
% clipped near target, off-scale init transient annotated on-plot)
\includegraphics[width=0.85\linewidth]{torus_nsig.png}
\caption{Torus phantom check (main paper \S5.1): significant-H1 count over
training (seed means; dashed line: true $\beta_1{=}2$). After the brief
resampling transient near step 100, every loss arm sits exactly on two
significant loops --- zero phantom handles --- while the repulsion control
(C2) collapses a loop by step ${\sim}800$ and never recovers it.}
\label{fig:nsig}
\end{figure}

\paragraph{The stage-3a gate (main paper \S5).} Fig.~\ref{fig:toys} shows
three of the four repair probes run before integration: the loss alone
(no photometric term, Adam directly on point positions) must repair
defective clouds. The predicted location-blind pathology of recruitment
did not materialize, for a mechanism reason worth stating: tearing a
2-manifold shell cannot grow a component bar past the local detour scale
(the tug is abandoned), while cutting the 1-D bridge raises the recruited
bar's death monotonically toward target --- the runaway concentrates on
the only profitable direction.
% src: figures/phase3_toy/t{1..4}.json (PHASE3_PLAN.md Appendix A)

\begin{figure}[t]
\centering
% src: figures/phase3_toy/t{1,2,4}.png (tests/test_topo_loss.py artifacts;
% title bands cropped for print, panel data untouched; factors in t*.json)
\includegraphics[width=\linewidth]{toy_t1.png}\\[2pt]
\includegraphics[width=\linewidth]{toy_t2.png}\\[2pt]
\includegraphics[width=\linewidth]{toy_t4.png}
\caption{The stage-3a gate: the loss alone (Adam on raw points, no
photometric term) repairs defective clouds; each row shows initial cloud,
final cloud, and bottleneck-to-target trajectory. \emph{Top} (T1):
punctured sphere --- the matched birth term closes the void ($8.8\times$).
\emph{Middle} (T2): spurious handle --- the diagonal term kills the extra
H1 bar. \emph{Bottom} (T4): bridged merge with no significant live H0 bar
--- recruitment severs the bridge ($450\times$, $\beta_0\,1{\to}2$), the
case where plain optimal matching provably has zero gradient. (T3,
mis-spaced components, $1229\times$: omitted for space.)}
\label{fig:toys}
\end{figure}

% Round-4 (2026-07-17): the decisive-shape trajectory figure moved here
% from main §5.1 (page budget for the round-4 teaser + qualitative
% matrix). Appended AFTER fig:toys so S2-S5 keep their numbers; this is
% Fig. S6. The full-width canvas also answers the review's Fig-1
% legibility note (many overlapping arms). Caption verbatim from main.
\begin{figure*}[t]
\centering
% src: topo3/report/{sphere,torus}_series.png (topo_loss_report.py; line =
% seed mean, band = seed min-max; underlying numbers in results.json)
\includegraphics[width=0.495\linewidth]{sphere_series.png}\hfill
\includegraphics[width=0.495\linewidth]{torus_series.png}
\caption{Topology error over training: bottleneck to target in the
discriminating dimension (line = seed mean, band = seed min--max).
\emph{Left}, sphere (H2): every loss arm (C1, C3, C5, C6, both $\rho$
variants, noise pair C7/C7h) converges to a fraction of baseline; the
norm-matched repulsion control C2 is \emph{destructive}, and its gentle
variant C2g still ends worse than baseline. \emph{Right}, torus (H1), the
class no prior won: same ordering, C5 best; noise arms descend shallower
but never approach baseline; the no-recruitment ablation C6 rides it.}
\label{fig:mainseries}
\end{figure*}

% Round-4 (2026-07-17): fig:tomyum moved here from main §5.2 (page
% budget); appended after fig:mainseries so it is Fig. S7. The pot's
% showcase render re-enters main inside the qualitative matrix. Caption
% verbatim except the cross-doc \ref -> plain "main paper §4".
\begin{figure}[t]
\centering
% src: figures/tomyum_pot_alu{,_cut}.png (scripts/make_kinkin_figure.py ---
% CPU ray-traced render of the certified shell _meshes/kinkin_src.ply;
% aluminium material per advisor request 2026-07-11; rendered name
% tom-yum pot = internal asset name kinkin)
\includegraphics[width=0.365\linewidth]{tomyum_pot_alu.png}\hfill
\includegraphics[width=0.365\linewidth]{tomyum_pot_alu_cut.png}
\caption{The \emph{tom-yum pot}: an artist-authored Blender model (raw
export: 9 open shells, 232 non-manifold junction edges) ingested and
certified to a single genus-3 shell, $\beta{=}(1,6,1)$ exact
(main paper \S4). \emph{Right}: cutaway --- the narrow flue mouth caps
at small $\alpha$, so the flue interior reads as an enclosed chamber: the
on-axis H2 void the floor rule selects at $M{=}8192$ (all H1 loops sit
below the floor).}
\label{fig:tomyum}
\end{figure}

% Round-4: the loop-class row of the qualitative matrix (main Fig. 1
% holds the two H2 rows under the page budget). Same renderer, badges,
% and camera discipline as Fig. 1; Fig. S8.
% src: scripts/make_matrix_figure.py (bob strip; seed-0 final_params.pt;
% badge = results.json nsig_final, identical across seeds)
\begin{figure*}[t]
\centering
\includegraphics[width=0.9\linewidth]{matrix_bob.png}
\caption{The loop-class row of the main paper's Fig.~1 matrix: bob
(genus 1, H1). The repulsion control's reconstruction reads one
significant loop against the true two ($\beta_1\,2{\to}1$, all seeds)
at $2\times$ worse Chamfer, while baseline and loss both read two ---
the loss at half the baseline's tail error (main Table~2).}
\label{fig:bobmatrix}
\end{figure*}
